\documentclass[10pt,a4paper]{amsart}
\usepackage[margin=19mm]{geometry}
\usepackage{amsmath,amssymb,mathtools}
\usepackage[scr=rsfs,cal=euler]{mathalfa}
\usepackage{xfrac}
\usepackage[unicode,hidelinks,bookmarks=false]{hyperref}
\newcommand{\ie}{i.e.\ }
\newcommand{\cf}{cf.\ }
\newcommand{\resp}{respectively\ }
\newcommand{\Subsection}{\subsection}
\providecommand{\nobreakdash}{\nobreak\hbox{-}}
\setlength{\emergencystretch}{2em}
\multlinegap0pt
\usepackage{amssymb}
\usepackage{cancel}
\usepackage{tikz}
\usepackage{float}
\usepackage{mathtools}
\newtheorem{theorem}{Theorem}
\def\EE{\mathbb{E}}
\newcommand{\PP}{\mathbb{P}}
\def\beq{\begin{equation}}
\def\dim{\mathrm{dim} }
\def\eeq{\end{equation}}
\title{Cohomology of moduli spaces of pointed curves}
\author{Jinwon Choi, Young-Hoon Kiem}
\date{Source: arXiv:2608.24052v1 (CC BY 4.0)}
\begin{document}
\maketitle

\noindent\emph{Source and licence.} Cohomology of moduli spaces of pointed curves. Version arXiv:2608.24052v1 (CC BY 4.0), distributed by arXiv under the Creative Commons Attribution 4.0 International licence, http://creativecommons.org/licenses/by/4.0/, as recorded in the arXiv licence metadata for this exact version and shown on the abstract page https://arxiv.org/abs/2608.24052v1. Full licence notice: \texttt{licenses/arxiv-2608.24052-CC-BY-4.0.txt}.\par\smallskip

\noindent\textit{Editorial scope.} Counted unit: the article's theorem 2\_1 (source0.tex lines 439--442). The article's complete original proof of it is delivered with it (source0.tex lines 922--933, one \texttt{proof} environment of 12 lines, which the article places away from the statement). No mathematics is written, repaired or re-derived: the statement and the proof are the article's own lines, in the article's own notation, and no retained line is substituted or re-flowed.  Retained uncounted as the source-local context the proof needs: lines 871--884.  Nothing was omitted from any retained span, so the package carries no omission record.  The retained lines cite 1 works, whose entries are transcribed verbatim from the article's own bibliography source in the e-print and delivered with short editorial labels recorded in the item's reference mapping.  No retained line contains a figure environment, an image-inclusion command or drawing code, so no image is delivered and the package declares no asset dependency.  Editorial formatting only, in addition: an AMS \texttt{amsart} preamble that restates the article's own declarations and macros used by the retained lines, namely the environments \texttt{theorem} on separate counters, together with the standard packages those lines need.  The retained environments print in this document as Theorem 1 in order of appearance, whereas the article prints the same items with the numbering of its own full document.  The complete native source of the article as distributed in the arXiv e-print for this exact version is bundled unchanged as \texttt{sources/208445/source0.tex}.
\begin{theorem}\cite[Theorem 1]{HwangLLT}\label{LLT}
  Let $\{\Omega_n\}$ be a sequence of integer-valued random variables with moment generating functions $M_n(s)=\EE[e^{\Omega_n s}]$. Suppose there exist $\delta>0$ and analytic functions $u(s)$, $v(s)$ such that uniformly for $|s|<\delta$,
  \beq\label{eq:mgf} M_n(s) = e^{n u(s) +v(s)}\left(1+ O(n^{-1}) \right), \quad n\to \infty, \eeq
  satisfying the following conditions:
  \begin{enumerate}
    \item $u(s)$ and $v(s)$ are independent of $n$ and analytic for $|s|\le \delta$ with $u''(0)\ne 0$.
    \item \emph{(Aperiodicity condition)} For a fixed $0<\epsilon<\delta$, there exists $c>0$ such that
    \[ \left| \frac{M_n(r+it)}{M_n(r)}\right|=O(e^{-cn})\]
    uniformly for $|r|\le \delta$ and $\varepsilon \le |t|\le \pi$.
  \end{enumerate}
  Then, letting $\mu_n = n u'(0) +v'(0) $ and $\sigma_n^2 = n u''(0)$, for any integer $m$ in the central region such that $x =\frac{m-\mu_n}{\sigma_n}=O(1)$, the probabilities satisfy asymptotically
  \[ \PP(\Omega_n =m) = \frac{e^{-x^2/2}}{\sqrt{2 \pi \sigma_n^2}}\left( 1 + \frac{\Pi_1(x) }{\sqrt{n}} + \frac{\Pi_2(x) }{n}+ O\left( n^{-\frac{3}{2}}\right) \right),\]
  where $\Pi_1(x) $ and $\Pi_2(x) $ are polynomials.
\end{theorem}

\begin{theorem}\label{2_1}
 For any connected smooth projective curve $C$, let $\beta_{n,k}= \dim H^k(C[n])$.
 The subsequences of even Betti numbers $\{\beta_{n,2i}\}$ and odd Betti numbers $\{\beta_{n,2i+1}\}$ are asymptotically 3-ULC but not 4-ULC in the central range $|i - \frac{n}{2}|=O(\sqrt{n})$.
\end{theorem}

\begin{proof}[Proof of Theorem \ref{2_1}]
  Let $\tilde{p}_{n,i}= [w^i]\tilde{P}_n(w) $. Note that $\tilde{p}_{n,i}= \beta_{n,2i} / \sum_j \beta_{n,2j}$ is the normalized probability corresponding to the even Betti numbers. Hence, to show the sequence $\{\tilde{p}_{n,i}\}$ is asymptotically $r$-ULC, it suffices to show that the second discrete difference of the logarithmic ratio
  \[\Delta^2_i \ln \left(\frac{\tilde{p}_{n,i}}{\binom{n}{i}^r}\right) =\ln \left(\frac{\tilde{p}_{n,i-1}}{\binom{n}{i-1}^r}\right)-2 \ln \left(\frac{\tilde{p}_{n,i}}{\binom{n}{i}^r}\right)+\ln \left(\frac{\tilde{p}_{n,i+1}}{\binom{n}{i+1}^r}\right) \]
  is asymptotically strictly negative.
  By Theorem \ref{LLT}, for $i$ in the central region $|i - \frac{n}{2}|=O(\sigma_n)=O(\sqrt{n})$,
  \[ \Delta^2_i \ln \tilde{p}_{n,i} = - \frac{1}{\sigma_n^2}+ O\left( n^{-\frac{3}{2}}\right), \]
  and since $\Delta^2_i \ln {\binom{n}{i}^r}= - \frac{4r}{n}+ O\left( n^{-2}\right)$, the required inequality evaluates to
  \[\Delta^2_i \ln \left(\frac{\tilde{p}_{n,i}}{\binom{n}{i}^r}\right) = - \frac{1}{\sigma_n^2}+\frac{4r}{n}+O\left( n^{-\frac{3}{2}}\right)<0.\]
  Because the error term $O\left( n^{-\frac{3}{2}}\right)$ decays asymptotically faster than the main term of order $n^{-1}$, the inequality holds for sufficiently large $n$ if and only if $r< \frac{n}{4\sigma_n^2}$. Evaluating the threshold ratio, we get
  \[ \frac{n}{4\sigma_n^2} = \frac{1}{4 u''(0)}= \frac{6(e-2)}{4 (3-e)}\approx 3.82. \]
  Therefore the subsequence is asymptotically 3-ULC but not 4-ULC. The proof for the subsequence of the odd Betti numbers is parallel.
\end{proof}

\begin{thebibliography}{99}
\bibitem[HwangL]{HwangLLT} H.-K. Hwang. {\em Large deviations of combinatorial distributions. II. Local limit theorems}. Ann. Appl. Probab. 8(1): 163-181 (1998).
\end{thebibliography}
\end{document}
